\chapter{Wholesaling du flux de détail}\label{ch:hf:retail-wholesaling}

Un ordre de détail d'achat de cinquante actions, passé depuis un téléphone, est acheminé vers un wholesaler, qui l'exécute à un prix légèrement meilleur que le meilleur cours vendeur sur n'importe quelle bourse,
garde la transaction et verse au courtier une rémunération pour ce privilège. Dans le modèle de ce chapitre, un wholesaler qui consacre aux ordres de détail 30~\% de la demi-fourchette en amélioration de prix
et verse un dixième de centime par action pour ces ordres conserve 0.44 centime par action après couverture~; le trader de détail paie une fourchette effective égale à 70~\% de la fourchette affichée.
Les mêmes conditions appliquées aux ordres institutionnels feraient perdre 0.94 centime par action~: l'activité est rentable grâce à la
segmentation.

\section{Économie du flux de détail}

Le livre~1, chapitre~10 décrit cette structure~: les courtiers de détail acheminent les ordres au marché de leurs clients vers des wholesalers, des teneurs de marché hors bourse,
qui les exécutent pour compte propre, en général au niveau des meilleurs cours acheteur et vendeur nationaux (NBBO) ou à l'intérieur de ceux-ci, et versent aux courtiers un paiement pour flux d'ordres (PFOF). Le revenu du wholesaler
par action comporte cinq composantes~:
\[
\text{net}=\underbrace{h(1-\pi)}_{\text{gain capté}}-\underbrace{a}_{\text{mark-out}}-\underbrace{p}_{\text{paiement pour flux}}
-\underbrace{c_H}_{\text{couverture}}+\underbrace{I}_{\text{P\&L d'inventaire}},
\]
avec $h$ la demi-fourchette, $\pi$ la part de celle-ci rendue en amélioration de prix, $a$ l'information contenue dans l'ordre, $p$ le paiement pour flux et $c_H$ le
coût de la couverture de ce que le wholesaler ne veut pas détenir. Tout repose sur $a$~: les ordres de détail sont petits, nombreux et, sur la minute suivante, presque
dénués d'information~; les ordres institutionnels ne le sont pas.

\begin{definition}[Segmentation des flux]\label{def:hf:retail-wholesaling:segmentation}
La \emph{segmentation des flux}\index{segmentation des flux} consiste à séparer les flux d'ordres selon leur origine (détail, institutionnel, pour compte propre) afin que chaque segment
soit exécuté à des prix fixés par sa propre information et son propre coût attendus, plutôt qu'à un prix unique fixé par le mélange~; c'est ce qui permet à un wholesaler
d'améliorer les cotations que le flux informé a élargies.
\end{definition}

\begin{omfigure}
\begin{tikzpicture}[font=\footnotesize,>=stealth,box/.style={draw,rounded corners,fill=omThm!8,minimum width=2.2cm,minimum height=0.75cm,align=center}]
\node[box] (r) at (0,0) {trader de détail\\(achat de 50)};
\node[box] (b) at (3.2,0) {courtier};
\node[box,fill=omDef!10] (w) at (6.8,0) {wholesaler};
\node[box] (x) at (11.0,1.1) {bourses\\(NBBO)};
\node[box] (h) at (11.0,-1.1) {couverture sur\\une bourse};
\draw[->] (r) -- node[above] {ordre} (b);
\draw[->] (b) -- node[above] {routage} (w);
\draw[->,dashed] (w.north) to[bend right=25] node[above] {paiement pour flux} (b.north);
\draw[->,dotted] (x.west) -- (w.north east);
\node[font=\scriptsize,anchor=south] at (8.9,0.95) {cotation};
\draw[->] (w.south east) -- (h.west);
\node[font=\scriptsize,anchor=north,align=center] at (8.7,-0.95) {inventaire\\excédentaire};
\draw[->,dashed] (w.south) to[bend left=25] node[below] {exécution au NBBO moins l'amélioration} (r.south);
\end{tikzpicture}

\omcaption{Le parcours d'un ordre de détail~: le courtier l'achemine vers un wholesaler, qui l'exécute pour compte propre au meilleur cours acheteur ou vendeur national
amélioré d'une partie de la demi-fourchette, verse au courtier un paiement pour flux d'ordres, et couvre sur une bourse l'inventaire qu'il ne souhaite pas détenir. Source~: livre~1, chapitre~10~; ce
chapitre.}\label{fig:hf:retail-wholesaling:path}
\end{omfigure}

\section{Segmentation et tarification}

La journée du chapitre (\cref{lst:hf:retail-wholesaling:flow})~: 20\,000 ordres de détail arrivant environ une fois par seconde, des tailles lognormales autour d'une médiane de 50
actions, 0.1 centime par action d'information dans le sens de l'ordre, des sens d'ordre dictés par un sentiment grégaire qui pousse les traders de détail à acheter (ou à vendre) par vagues~;
l'action est cotée avec une fourchette de 2 à 3 centimes. Pour comparaison, mille ordres institutionnels autour d'une médiane de 2\,000 actions, avec 1.5 centime d'information.
Le wholesaler exécute chaque ordre à la cotation améliorée de $\pi$ de la demi-fourchette, verse 0.1 centime par action et détient jusqu'à 10\,000 actions d'inventaire
et couvre l'excédent sur une bourse à la demi-fourchette, plus des frais de 0.3 centime (\cref{lst:hf:retail-wholesaling:wholesale}).

\begin{omfigure}
\begin{tikzpicture}
\begin{axis}[width=11cm,height=5.4cm,xlabel={\small amélioration de prix (\% de la demi-fourchette)},ylabel={\small centimes par action},xmin=0,xmax=80,
  ymin=-1.7,ymax=1.0,tick label style={font=\footnotesize},grid=major,grid style={gray!25},table/col sep=comma,
  legend style={font=\scriptsize,at={(0.5,-0.3)},anchor=north},legend columns=3]
\addplot[omThm,thick,mark=*,mark size=1.3pt] table[x=pi,y=retail]{figdata/hft/23-retail-wholesaling/pi.csv};
\addplot[omDef,thick,dashed,mark=square*,mark size=1.3pt] table[x=pi,y=inst]{figdata/hft/23-retail-wholesaling/pi.csv};
\addplot[gray,thick,dotted,mark=triangle*,mark size=1.3pt] table[x=pi,y=improvement]{figdata/hft/23-retail-wholesaling/pi.csv};
\draw[black,thin] (axis cs:0,0) -- (axis cs:80,0);
\legend{{wholesaler, flux de détail},{wholesaler, flux institutionnel},{amélioration pour le trader}}
\end{axis}
\end{tikzpicture}

\omcaption{Le revenu net par action du wholesaler avant P\&L d'inventaire (capture diminuée du mark-out, du paiement pour flux et de la couverture) sur le flux de détail et sur le flux institutionnel
et l'amélioration de prix reçue par le trader de détail, en fonction de la part de la demi-fourchette rendue~; paiement pour flux 0.1 centime par action,
limite d'inventaire 10\,000 actions. Données~: \texttt{hf\_wholesale.by\_pi}.}\label{fig:hf:retail-wholesaling:pi}
\end{omfigure}

Sans amélioration, le wholesaler conserve 0.85 centime par action sur le flux de détail et perd 0.57 sur le flux institutionnel
(\cref{fig:hf:retail-wholesaling:pi}). Chaque tranche de 10~\% de la demi-fourchette rendue lui coûte 0.125 centime~; l'activité de détail atteint le seuil de rentabilité à 68~\% et
n'atteint jamais celui de l'activité institutionnelle. Cet écart représente la valeur de la segmentation, et c'est lui qui finance à la fois l'amélioration et le paiement pour flux.

\section{Amélioration de prix et sa mesure}

\begin{definition}[Taux d'internalisation]\label{def:hf:retail-wholesaling:rate}
Le \emph{taux d'internalisation}\index{taux d'internalisation} est la part des exécutions d'un teneur de marché (en actions ou en ordres) qu'il conserve
face à son propre inventaire, plutôt que de la compenser sur une bourse ou une autre plateforme dans un court intervalle de temps.
\end{definition}

L'amélioration de prix se mesure par rapport aux meilleurs cours acheteur et vendeur nationaux à l'arrivée de l'ordre. Les statistiques de la Rule 605 du livre~1, chapitre~10 la rapportent
par centre de marché~: la fourchette effective (le double de la distance entre le mid et l'exécution), son rapport à la fourchette affichée (E/Q) et l'amélioration par
action. Dans le modèle, à 30~\% d'amélioration, le trader de détail fait face à une fourchette affichée de 2.5 centimes et à une fourchette effective de 1.75~: E/Q vaut 0.70, et
l'amélioration est de 0.375 centime par action.

L'amélioration par rapport à la cotation n'équivaut pas à un bon prix. Schwarz, Barber, Huang, Jorion et Odean ont placé 85\,000 ordres au marché simultanés
par six comptes auprès de cinq courtiers~: le coût aller-retour moyen au niveau des comptes variait de 0.07~\% à 0.46~\%, hors commissions, parce que les
wholesalers proposaient systématiquement des prix différents pour les mêmes transactions selon les courtiers~; la variation du paiement pour flux d'ordres n'expliquait pas
cet écart. Le \omterm{def:hf:retail-wholesaling:rate}{taux d'internalisation} est ici de 87~\%~: les 13~\% du volume de détail que le wholesaler couvre sont le flux arrivé par vagues.

\begin{dated}{2026-09}{Les règles sur l'exécution des ordres de détail}\label{dat:hf:retail-wholesaling:rules}
Dans son rapport annuel pour 2025, Virtu Financial indique qu'en juin 2025 la SEC a retiré ses projets en cours pour l'Order Competition Rule
(Rule 615 proposée de la Regulation NMS), pour la Regulation Best Execution et pour une restriction de la tarification à paliers des bourses fondée sur le volume, et que la
date de mise en conformité des modifications de la Rule 605 de la Regulation NMS, fixée auparavant au 15 décembre 2025 environ, a été reportée au 1\up{er} août 2026. Le 3 juillet
2012, la SEC a approuvé à titre pilote le Retail Liquidity Program de la New York Stock Exchange, avec des ordres d'amélioration de prix de détail non affichés,
fixés à un prix meilleur que le meilleur cours acheteur ou vendeur protégé d'au moins 0.001~\$ par action.
\end{dated}

\section{Inventaire issu d'un flux à sens unique}

Le flux de détail se compense sur une journée mais pas sur quelques minutes~: lorsque de nombreux traders de détail achètent le même titre en même temps, le wholesaler est vendeur de l'ensemble. La
limite d'inventaire arbitre entre le coût de couverture et le risque (\cref{fig:hf:retail-wholesaling:limits}). Avec une limite de 1\,000 actions, le wholesaler couvre 0.47
centime par action et obtient 0.20 avec un écart-type quotidien de 0.06~; à 10\,000, il couvre 0.20 et obtient 0.45 avec un écart-type de 0.38~; à 50\,000
sa couverture disparaît presque, mais son résultat quotidien est dominé par l'inventaire, avec un écart-type de 1.32 centime par action. Le grégarisme détermine la part
qui doit réellement être couverte~: sans lui, la même limite couvre 0.9~\% du volume~; avec la persistance du sentiment du modèle, 12.9~\%.

\begin{omfigure}
\begin{tikzpicture}
\begin{axis}[width=11cm,height=5.2cm,xlabel={\small taille de l'inventaire limite (milliers d'actions)},ylabel={\small cents nets par action (moyenne $\pm$ écart-type)},
  xmode=log,log basis x=10,xtick={1,2,5,10,20,50},xticklabels={1,2,5,10,20,50},xmin=0.8,xmax=65,ymin=-1.4,ymax=1.9,
  tick label style={font=\footnotesize},grid=major,grid style={gray!25},table/col sep=comma,
  legend style={font=\scriptsize,at={(0.03,0.97)},anchor=north west}]
\addplot[omThm,thick,mark=*,mark size=1.4pt,error bars/.cd,y dir=both,y explicit] table[x=limit,y=mean,y error=sd]{figdata/hft/23-retail-wholesaling/limits.csv};
\addplot[omDef,thick,dashed,mark=square*,mark size=1.3pt] table[x=limit,y=hedge]{figdata/hft/23-retail-wholesaling/limits.csv};
\legend{{net par action, moyenne $\pm$ écart-type sur 20 jours},coût de couverture par action}
\end{axis}
\end{tikzpicture}

\omcaption{Le résultat quotidien net par action du wholesaler (moyenne et écart-type sur vingt journées simulées de flux de détail) et son coût de couverture par action,
en fonction de l'inventaire détenu avant couverture sur une bourse~; amélioration de prix de 30~\%, 0.1 cent par action payé pour le flux. Données~:
\texttt{hf\_wholesale.limits}.}\label{fig:hf:retail-wholesaling:limits}
\end{omfigure}

\section{Règles et débat}

Le débat sur le paiement pour flux d'ordres (livre~1, chapitre~10) porte sur la question de savoir qui capte la valeur de la segmentation~: le wholesaler, le courtier par ses paiements,
ou le trader de détail par l'amélioration de prix. L'arithmétique du chapitre montre la marge disponible~: 1.25 cent de demi-fourchette, moins 0.1 d'information et 0.2
de couverture, laisse 0.95 cent par action à répartir~; le modèle attribue 0.375 au trader et 0.1 au courtier, et le wholesaler conserve 0.475 (0.44
après le P\&L de son inventaire). La Order Competition Rule, retirée, aurait envoyé de nombreux ordres de détail vers des enchères
où d'autres traders auraient pu se les disputer~; les \omterm{def:hf:retail-wholesaling:rlp}{programmes de liquidité de détail} des bourses, celui de la NYSE approuvé en 2012, permettent aux membres de la bourse d'offrir
aux ordres de détail une amélioration de prix inférieure au cent, directement sur la bourse elle-même.

\begin{definition}[Programme de liquidité de détail]\label{def:hf:retail-wholesaling:rlp}
Un \emph{programme de liquidité de détail}\index{programme de liquidité de détail} (retail liquidity programme) est un mécanisme de bourse dans lequel les ordres identifiés comme provenant du détail peuvent être exécutés
contre des ordres d'amélioration de prix non affichés que seuls les ordres de détail peuvent prendre, cotés à l'intérieur du meilleur cours acheteur ou vendeur d'une fraction de cent.
\end{definition}

\section{Fiches de stratégie}

\begin{strategyfile}{Internalisation des actions de détail}\label{strat:hf:retail-wholesaling:equity}
\sfield{Qui paie, et pourquoi} Les traders de détail, par la part de la fourchette non rétrocédée, pour l'immédiateté et un prix égal ou meilleur que le cours affiché.
\sfield{Instruments et plateformes} Actions américaines~; flux routés par les courtiers~; bourses pour la couverture.
\sfield{Signal} Le cours affiché à l'arrivée~; le segment du flux et l'information attendue par courtier et type d'ordre.
\sfield{Dimensionnement et exécution} Exécuter chaque ordre au cours affiché amélioré de la part convenue~; conserver jusqu'à la limite d'inventaire~; couvrir au-delà.
\sfield{Coûts} Paiement pour flux, couverture (0.20 cent par action pour une limite de 10\,000 actions dans le modèle), mark-outs, risque de l'inventaire.
\sfield{Comment elle s'éteint} Un flux moins de détail qu'il n'y paraît, et la réglementation du routage~; l'épisode cité est la constatation, par Schwarz, Barber, Huang, Jorion et
Odean, que le même ordre obtenait des prix différents selon les courtiers.
\sfield{Horizon, capacité, infrastructure} Millisecondes pour l'exécution, minutes pour la couverture~; les accords de flux d'ordres des courtiers.
\sfield{Backtest honnête} Le cours affiché à l'arrivée, le mark-out ultérieur de l'ordre par courtier, les journées de mimétisme.
\sfield{Sources} Schwarz et al.\ (2025)~; livre~1, chapitre~10~; le présent chapitre.
\end{strategyfile}

\begin{strategyfile}{Flux d'options de détail par enchères d'amélioration de prix}\label{strat:hf:retail-wholesaling:options}
\sfield{Qui paie, et pourquoi} Les traders d'options de détail, dont les ordres portent des fourchettes larges et peu d'information.
\sfield{Instruments et plateformes} Options cotées~; enchères d'amélioration de prix des bourses (livre~1, chapitre~24), où la filiale du wholesaler peut
garantir l'ordre.
\sfield{Signal} Les fair values des options (chapitre~19) et le segment de l'ordre.
\sfield{Dimensionnement et exécution} Garantir et améliorer l'ordre dans l'enchère~; couvrir le delta sur l'action.
\sfield{Coûts} Paiement pour flux (plus élevé pour les options que pour les actions), règles d'attribution de l'enchère, couverture.
\sfield{Comment elle s'éteint} Des règles d'enchère qui donnent davantage de l'attribution aux concurrents~; les résultats d'Ernst et Spatt selon lesquels l'internalisation des options est
imparfaitement concurrentielle, ce qui justifie les paiements.
\sfield{Horizon, capacité, infrastructure} Millisecondes~; connectivité aux enchères sur chaque bourse d'options.
\sfield{Backtest honnête} L'attribution de l'enchère de chaque bourse à la date considérée.
\sfield{Sources} Ernst et Spatt (2022, document de travail)~; livre~1, chapitre~24.
\end{strategyfile}

\begin{strategyfile}{Programmes de liquidité de détail des bourses}\label{strat:hf:retail-wholesaling:rlp}
\sfield{Qui paie, et pourquoi} Les ordres de détail que les courtiers envoient au programme de détail d'une bourse.
\sfield{Instruments et plateformes} Programmes de détail des bourses~; ordres d'amélioration de prix non affichés.
\sfield{Signal} La même segmentation, sur la bourse~: l'indicateur de détail.
\sfield{Dimensionnement et exécution} Laisser au repos des ordres d'amélioration inférieure au cent que seuls les ordres de détail peuvent prendre~; couvrir sur le carnet affiché.
\sfield{Coûts} Frais de la bourse~; l'amélioration accordée.
\sfield{Comment elle s'éteint} Des courtiers qui routent le flux de détail vers des wholesalers à la place~; l'épisode cité est le programme de la NYSE, approuvé en juillet 2012 à titre pilote.
\sfield{Horizon, capacité, infrastructure} Millisecondes~; adhésion à la bourse.
\sfield{Backtest honnête} Seuls les ordres de détail qui ont effectivement atteint le programme.
\sfield{Sources} Décision de la SEC approuvant le \omterm{def:hf:retail-wholesaling:rlp}{programme de liquidité de détail} de la NYSE (2012).
\end{strategyfile}

\section{Tutoriel~: cinquante actions depuis un téléphone}\label{tut:hf:retail-wholesaling:tut}

\begin{tutorial}
\textbf{Objectif.} Décomposer le résultat d'un wholesaler sur un flux segmenté et déterminer combien d'amélioration de prix et combien d'inventaire il peut se permettre.
\textbf{État final~:} les \cref{fig:hf:retail-wholesaling:pi,fig:hf:retail-wholesaling:limits} et les statistiques de la Rule 605.
\begin{tutsteps}
\item \textbf{Le flux}~: le détail avec mimétisme, l'institutionnel avec information.
\omcode{code/firm/wholesale/firm_wholesale.py}{35}{49}{Tailles des ordres de détail, sentiment de mimétisme qui détermine leur sens, et leur information~; ordres institutionnels à titre de comparaison.}\label{lst:hf:retail-wholesaling:flow}
\item \textbf{Le wholesaler}~: capture, mark-out, paiement, couverture au-delà de la limite, P\&L d'inventaire.
\omcode{code/firm/wholesale/firm_wholesale.py}{60}{79}{Chaque ordre exécuté au cours amélioré~; inventaire couvert au-delà de sa limite~; les composantes par action.}\label{lst:hf:retail-wholesaling:wholesale}
\item \textbf{Balayer} l'amélioration et la limite (\texttt{hf\_wholesale.by\_pi}, \texttt{hf\_wholesale.limits})~; les statistiques de la Rule 605
(\texttt{firm.wholesale.rule605}).
\end{tutsteps}
\textbf{Pour aller plus loin.} Faire varier le paiement selon le courtier et observer qui obtient l'amélioration~; ajouter un momentum de prix aux journées de mimétisme~; router une partie du
flux vers un \omterm{def:hf:retail-wholesaling:rlp}{programme de liquidité de détail}.
\end{tutorial}

\section{Construction~: le wholesaler}\label{bld:hf:retail-wholesaling:wholesale}

\begin{build}
\bfield{Rôle} Simuler un flux segmenté, l'exécuter avec amélioration de prix, couvrir le résiduel, et présenter les composantes du résultat ainsi que des
statistiques de type Rule 605.
\bfield{Interface} \texttt{Flow(kind, n, seed, rate, herd, info, half\_spread, spread\_jitter)},
\texttt{wholesale(flow, pi\_share, pfof, limit, hedge\_fee, sigma\_day, seed)}, \texttt{rule605(flow, pi\_share)}.
\bfield{Règles} Cents par action~; amélioration exprimée en part de la demi-fourchette~; couverture à la demi-fourchette plus des frais.
\bfield{Tests d'acceptation} \texttt{code/firm/wholesale/tests/}~: les ordres de détail plus petits et moins informés que les ordres institutionnels~; la décomposition à
la main sans bruit ni limite~; davantage d'amélioration abaisse le net~; une limite plus serrée couvre davantage~; E/Q égale un moins la part d'amélioration.
\bfield{Pour aller plus loin} Paiements et amélioration propres à chaque courtier~; impact de prix du mimétisme~; flux d'options et enchères.
\end{build}

\begin{omsources}
\item C.~Schwarz, B.~M.~Barber, X.~Huang, P.~Jorion, T.~Odean, The ``actual retail price'' of equity trades, \emph{Journal of Finance} 80(5), 2025,
2507--2541.
\item T.~Ernst, C.~S.~Spatt, Payment for order flow and option internalization, document de travail, 2022.
\item Virtu Financial, rapport annuel sur le Form 10-K pour 2025.
\item US Securities and Exchange Commission, Release No.\ 34-67347 (NYSE Retail Liquidity Program), 3 juillet 2012.
\end{omsources}

\section{Exercices}

\begin{exercise}[$\star$]\label{exo:hf:retail-wholesaling:1}
Le cours affiché est 50.00--50.02 et le grossiste accorde 30~\% de la demi-fourchette. À quel prix un ordre d'achat de détail est-il exécuté, et quelle est sa
fourchette effective~?
\end{exercise}

\begin{exercise}[$\star$]\label{exo:hf:retail-wholesaling:2}
Avec une demi-fourchette de 1.25 cent, un mark-out de 0.1, un paiement de 0.1 et une couverture de 0.2, à partir de quelle part d'amélioration le grossiste atteint-il le seuil de rentabilité~?
\end{exercise}

\begin{exercise}[$\star$]\label{exo:hf:retail-wholesaling:3}
Quel ratio E/Q donne une amélioration de 30~\%, et pourquoi est-il indépendant de la fourchette~?
\end{exercise}

\begin{exercise}[$\star\star$]\label{exo:hf:retail-wholesaling:4}
Pourquoi le mimétisme conduit-il un grossiste à couvrir quatorze fois plus de volume à la même limite~?
\end{exercise}

\begin{exercise}[$\star\star$]\label{exo:hf:retail-wholesaling:5}
Pourquoi le grossiste ne pourrait-il pas offrir les mêmes conditions aux flux institutionnels~?
\end{exercise}

\begin{exercise}[$\star\star$]\label{exo:hf:retail-wholesaling:6}
Schwarz et ses coauteurs ont constaté une dispersion des prix entre courtiers que le paiement pour flux d'ordres n'explique pas. Qu'est-ce qui pourrait l'expliquer d'autre~?
\end{exercise}

\begin{exercise}[$\star\star\star$]\label{exo:hf:retail-wholesaling:7}
\emph{Programmation.} Relancer \texttt{hf\_wholesale.limits}~: quelle limite maximise le net moyen, et quelle limite maximise le net moyen diminué de l'écart type~?
\end{exercise}

\begin{exercise}[$\star\star\star$]\label{exo:hf:retail-wholesaling:8}
\emph{Trouver la faille.} «~Notre E/Q vaut 0.70, donc nos clients récupèrent 30~\% de la fourchette~: la meilleure exécution du marché.~»
\end{exercise}

\section{Problème~: cinquante actions depuis un téléphone}

\begin{problem}[{Problème du week-end --- cinquante actions depuis un téléphone}]\label{pb:hf:retail-wholesaling:1}
Un grossiste fixe ses conditions pour le flux de détail d'un courtier.

\medskip\noindent\textbf{Partie I --- L'activité.}
\begin{enumerate}
\item Écrire le net par action du grossiste et nommer ses composantes.
\item Définir la \omterm{def:hf:retail-wholesaling:segmentation}{segmentation des flux}.
\item Décrire les flux de détail et institutionnels du modèle.
\item À quoi la segmentation rapporte-t-elle~?
\end{enumerate}

\medskip\noindent\textbf{Partie II --- Amélioration.}
\begin{enumerate}[resume]
\item Donner le net par action à 0, 30~\% et 70~\% d'amélioration, pour le flux de détail et le flux institutionnel.
\item Où l'activité de détail atteint-elle le seuil de rentabilité~?
\item Définir le \omterm{def:hf:retail-wholesaling:rate}{taux d'internalisation} et donner celui du modèle.
\item Donner les statistiques de la Rule 605 à 30~\%.
\end{enumerate}

\medskip\noindent\textbf{Partie III --- Inventaire et règles.}
\begin{enumerate}[resume]
\item Que met en balance la limite d'inventaire~? Donner trois points.
\item Que fait le mimétisme à la couverture~?
\item Résumer l'encadré daté.
\item Définir un \omterm{def:hf:retail-wholesaling:rlp}{programme de liquidité de détail} (retail liquidity programme).
\end{enumerate}

\medskip\noindent\textbf{Partie IV --- Le verdict.}
\begin{enumerate}[resume]
\item Énoncer le \emph{résultat nommé}~: le net par action du grossiste, une fois déduits le paiement pour flux d'ordres et la couverture, et l'amélioration de prix du trader de détail
par rapport à la fourchette affichée.
\item Qui capte la valeur de la segmentation~?
\item Qu'ont constaté Schwarz et ses coauteurs, et qu'implique cela pour les informations publiées~?
\item Qu'aurait changé une règle de mise en concurrence des ordres~?
\item Quelle fiche de stratégie est la plus exposée à la réglementation~?
\item Comment détecter un flux institutionnel déguisé en flux de détail~?
\item Que doit demander un courtier à un grossiste~?
\item En une phrase~: que achète un grossiste~?
\end{enumerate}
\end{problem}

\section{Questions d'entretien}

\begin{interviewq}[$\star$ \iqroles{trader}]\label{iq:hf:retail-wholesaling:1}
Pourquoi un grossiste peut-il exécuter des ordres de détail à l'intérieur de la cotation alors que les teneurs de marché des bourses ne le peuvent pas~?
\end{interviewq}

\begin{interviewq}[$\star\star$ \iqroles{chercheur}]\label{iq:hf:retail-wholesaling:2}
Comment mesureriez-vous l'information contenue dans le flux d'ordres d'un courtier~?
\end{interviewq}

\begin{interviewq}[$\star\star$ \iqroles{développeur}]\label{iq:hf:retail-wholesaling:3}
Concevriez-vous le système qui fixe le prix d'un ordre de détail en 100 microsecondes à partir de la cotation, des conditions du courtier et de l'inventaire~?
\end{interviewq}

\begin{interviewq}[$\star\star$ \iqroles{risque}]\label{iq:hf:retail-wholesaling:4}
Une journée meme stock envoie une vague d'ordres d'achat de détail sur un même titre. Quelles limites s'appliquent, et que deviennent les cotations du grossiste~?
\end{interviewq}

\begin{interviewq}[$\star\star$ \iqroles{trader}]\label{iq:hf:retail-wholesaling:5}
Un courtier demande un paiement pour flux d'ordres plus élevé. Comment décidez-vous~?
\end{interviewq}

\begin{interviewq}[$\star\star\star$ \iqroles{chercheur}]\label{iq:hf:retail-wholesaling:6}
Le sentiment de mimétisme suit un AR(1) de persistance $\rho$ par ordre~; comment la variance du déséquilibre net des ordres sur $N$ ordres croît-elle
avec $\rho$~?
\end{interviewq}
